\paragraph{The smooth model and its Jacobian.}
The normalization \(\nu:\widetilde C\to C\) has genus 43. Each of the
two geometric singularities has two smooth branches with intersection
multiplicity three. The branches at either point are already defined
over its residue field \(K=\mathbb F_{25}\).

To see this, write \(s^2+s+1=0\) and work at
\((a,b,c,d)=(0,s,0,1)\) in the chart \(m=1\). Set \(t=d-1\), and let
\[
 M=\det\frac{\partial(q_2,q_4)}{\partial(a,b)},\qquad
 G=\det\frac{\partial(q_2,q_4,q_3)}{\partial(a,b,c)}.
\]
At this point \(M=-2s-1\ne0\). The formal implicit function theorem
solves \(q_2=q_4=0\) for \(a,b-s\) in \(K[[c,t]]\). Write \(F(c,t)\)
for \(q_3\) on this smooth surface. The chain rule gives \(G=M F_c\)
on the surface. The determinant
\(\det\bigl(\partial(q_2,q_4,G)/\partial(a,b,c)\bigr)\) at the point
is \(-1\). In particular \(F_{cc}(0,0)=-1/M^2=2\), and there is a
unique formal critical arc \(c=c_0(t)\) with \(F_c=0\), \(c_0(0)=0\).

The following rows give the coefficients of \(t,\ldots,t^6\) on this
arc. Direct substitution in the original permanent equations gives
\(q_2=q_4=G=0\bmod t^7\) and \(q_3=-2t^6\bmod t^7\).
\[
\begin{array}{c|rrrrrr}
 & t&t^2&t^3&t^4&t^5&t^6\\ \hline
 a   &1&1&0&2&0&-2\\
 b-s &2&s-2&-2s-1&2-s&2-s&2-s\\
 c   &s&2-s&2-s&2s+1&-2s-1&s-2
\end{array}
\]
Thus after translation by \(c_0(t)\) the equation is
\(X^2U(X,t)-t^6 V(t)=0\) with \(U(0,0)=1\) and \(V(0)=2\).
The unit \(2U/V\) has a square root with constant term one, since
2 is invertible. The new coordinate \(Y=X\sqrt{2U/V}\) gives
\[
        \widehat{\mathcal O}_{C_K,(0,s,0,1)}
                   \simeq K[[Y,t]]/(Y^2-2t^6).
\]
Indeed \((2s+1)^2=2\). Evaluation along the branches
\(Y=\pm(2s+1)t^3\) identifies this ring with the pairs of series in
\(K[[t]]\times K[[t]]\) congruent modulo \(t^3\). Its normalization
is the whole product, and the quotient has length three. Completion
preserves this delta invariant \cite{stacksDeltaCompletion}. Frobenius
gives the same calculation at the conjugate singularity. By the
normalization exact sequence and the previously computed arithmetic
genus, \(g(\widetilde C)=49-3-3=43\). One can construct this model over \(\mathbb F_5\) by blowing up the
reduced singular locus three times, taking the new reduced singular
locus as center at each step. In the completed charts put
\(Y=tY_1,\ Y_1=tY_2,\ Y_2=tY_3\) successively. The strict transforms
satisfy \(Y_1^2=2t^4,\ Y_2^2=2t^2,\ Y_3^2=2\). The first two transforms
each have a unique singular point above the center. At the third step
the two tangent directions are finite and distinct. The centers are
Frobenius-stable and the morphisms are isomorphisms off those centers.
The resulting smooth projective geometrically connected curve is the
normalization.

The Jacobian has dimension 43. For each prime \(\ell\ne5\) its Tate
module has rank 86 \cite{stacksPicardCurves}. The computed point counts
on \(C\) are 18 over \(\mathbb F_5\) and 64 over \(\mathbb F_{25}\).
The two singular points and their four branch points are all defined
over \(\mathbb F_{25}\), with none rational over \(\mathbb F_5\).
The counts on \(\widetilde C\) are consequently 18 and 66. If \(\pi\)
is the fifth-power Frobenius on its Jacobian, the curve trace formula
gives \(\operatorname{Tr}(\pi)=-12\) and
\(\operatorname{Tr}(\pi^2)=-40\). Its zeta numerator
\(\det(1-T\pi\mid V_\ell\operatorname{Jac}(\widetilde C))\) has degree
86 and begins \(1+12T+92T^2\) \cite{longCurveFrobenius}.
For an ample line bundle \(L\) on the projective Jacobian over
\(\mathbb F_5\), let \(H\) be the numerical class of the symmetric
ample bundle \(L\otimes[-1]^*L\). Frobenius pulls back \(H\) to \(5H\),
so \(\pi^\dagger\pi=[5]\) for its Rosati adjoint. For rational \(a,b\)
the trace identity \(2gQ_H(u)=\operatorname{Tr}(u^\dagger u)\) and
\(\operatorname{Tr}(\pi^\dagger)=\operatorname{Tr}(\pi)\) give
\[
 Q_H(a\,1+b\pi)=a^2-\frac{12}{43}ab+5b^2,\qquad
 Q_H(\pi-1)=\frac{270}{43}.
\]
This computes the positive degree form on the span of the identity
and Frobenius \cite{longAbelianComparison}. Its discriminant is
\((12/43)^2-20<0\), so two rational linear combinations of these
endomorphisms agree exactly when their two coefficients agree. Applying its equality
test to the common-line condition still requires constructing two
endomorphisms and proving that their equality detects that incidence.
